\documentclass[../../main.tex]{subfiles}
\graphicspath{{../../image/}{../../figure/}}
\begin{document}

\section{稀疏波的传播}\label{section:ch11-稀疏波的传播}

我们已经熟悉真正非线性的去稳定作用: 压缩波波前变陡并最终破裂, 产生激波。然而, 真正非线性同时也通过诱导稀疏波波前的传播而施加正则化影响。值得注意的是, 这一效应纯粹是几何的, 与初值的正则性完全无关。

\begin{theorem}\label{theorem:ch11-11-2-1}
对任意可容许解 $u$,
\begin{align}
\frac{f'(u(y\pm,t))-f'(u(x\pm,t))}{y-x}\leqslant\frac{1}{t},\quad -\infty<x<y<\infty,\ 0<t<\infty.\label{align:ch11-11-2-1}
\end{align}
\end{theorem}

\begin{proof}
固定 $x,y$ 和 $t$, $x<y$, $t>0$。设 $\xi(\cdot)$ 与 $\zeta(\cdot)$ 分别表示从 $(x,t)$ 与 $(y,t)$ 出发的极大或极小后向特征线。由\cref{theorem:ch11-11-3}, $\xi(0)=x-tf'(u(x\pm,t))$, $\zeta(0)=y-tf'(u(y\pm,t))$。又由\cref{corollary:ch11-11-2}, $\xi(0)\leqslant\zeta(0)$。这立即推出 \eqref{align:ch11-11-2-1}。证毕。
\end{proof}

注意 \eqref{align:ch11-11-2-1} 建立了 $f'(u(\cdot,t))$ 的单侧 Lipschitz 条件, 其 Lipschitz 常数与初值无关。由\cref{chapter:ch06-标量守恒律L1理论}中给出的标量守恒律一般理论, 初值属于 $L^\infty(-\infty,\infty)$ 的\eqref{align:ch11-11-1-1}的可容许解, 可实现为一列初值在 $(-\infty,\infty)$ 上局部有界变差的解的几乎处处极限。因此, \eqref{align:ch11-11-2-1} 对初值仅属于 $L^\infty(-\infty,\infty)$ 的可容许解也应成立。显然, \eqref{align:ch11-11-2-1} 蕴含: 对固定 $t>0$, $f'(u(\cdot,t))$, 从而也 $u(\cdot,t)$, 在 $(-\infty,\infty)$ 的任意有界区间上有有界变差。于是我们证明了: 由于真正非线性, 解一般比其初值更光滑。


\begin{theorem}\label{theorem:ch11-11-2-2}
初值属于 $L^\infty(-\infty,\infty)$ 的\eqref{align:ch11-11-1-1}的可容许解, 在 $(-\infty,\infty)\times(0,\infty)$ 上属于 $BV_{\mathrm{loc}}$, 并满足单侧 Lipschitz 条件\eqref{align:ch11-11-2-1}。
\end{theorem}

\end{document}
